\section{Introduction}

Procurement set-asides for small firms are among the most widely deployed government-procurement instruments worldwide \citep{bosio2022}. They were designed on a textbook intuition: bringing more bidders to a protected segment should reduce procurement prices. The intuition has not survived contact with the data. A decade of empirical evidence places the price cost of bidder-exclusion rules in the ten-to-fifteen percent range \citep{marion2007, krasnokutskaya2011, athey2013, nakabayashi2013, hyytinen2018}. What price regressions cannot deliver is the channel: whether the markup runs through the change in \emph{who} can bid, or through the change in \emph{how} each remaining bidder bids. The welfare ranking of bidder exclusion against alternative preference instruments turns on which mechanism dominates, and the policy fix differs by an order of magnitude depending on the answer. If the markup is extensive, the prescription is to widen the protected roster. If it is intensive, no roster expansion will close the gap; the only lever is to shrink the rule. On São Paulo's medical-supplies procurement alone, the structural exercise of this paper prices the rule at R\$\,\welfAnnualBRLLo--\welfAnnualBRLHi~million (US\$\,\welfAnnualUSDLo--\welfAnnualUSDHi~million) per year of welfare loss---a magnitude this paper traces to a single mechanism.

A \policyCutoffMonth{} administrative reversal in S\~ao Paulo state, Brazil, supplies the natural experiment. The state extended an existing SME-only procurement rule to medical and hospital supplies---the last major exception---on the country's largest centralized state-government procurement platform, the Bolsa Eletr\^onica de Compras (BEC). The reversal was a legal opinion internal to the state prosecutor's office, defended on isonomy grounds and uncorrelated with prices, competition, or supplier complaints in the affected product group; the timing was bureaucratic. A difference-in-differences against \controlGroupCount{} never-treated product groups, reported as design validation in Appendix~\ref{app:did}, estimates that absolute winning prices rose \didLowerPct{}--\didUpperPct{} percent on the same window used by the structural exercise, with SME participation roughly doubling. The DiD coefficient fixes the average effect on observed log prices; it does not say which mechanism produced the markup or what the rule costs in welfare terms. The structural exercise of this paper recovers both, anchored to the same magnitude.

Three magnitudes appear in this paper and should not be confused. The DiD coefficient (\didLowerPct{}--\didUpperPct{} percent) is the average effect on observed log winning prices, item and month fixed effects partialled out. The structural counterfactual ($+\bneEffectTotalNp$ in non-pharma, $+\bneEffectTotalPh$ in pharma) is the simulated mean of the second-order statistic under each counterfactual entry profile, scaled by the reference price; in $p_{S_1}$ units it is three to four times the DiD coefficient because the two objects average over different primitives. The welfare cost (\welfLossPctLthirtyNp{}~percent of $p_{S_1}$ in non-pharma, \welfLossPctLthirtyPh{}~percent in pharma at $\lambda = \lambdaMain$) adds the marginal-cost-of-public-funds distortion on top of the allocative wedge. Sections~\ref{sec:results} and~\ref{sec:welfare} develop these in turn; Appendix~\ref{app:did} reports the DiD.

This paper provides that structure. I specify an asymmetric independent-private-values reading of the platform's Preg\~ao English-reverse auction format. Bidders draw costs from type-specific distributions $F_c^{\text{SME}}$ and $F_c^{\neg}$ and compete under a descending clock; losers exit at cost and the last firm supplies the contract. Drop-out prices of losers therefore point-identify $F_c^k$ under \citet{haile2003}, without inverting a bid function or imposing parametric structure on bidding. An auction-level common scale is shrunk out following \citet{krasnokutskaya2011}; winners are treated as left-censored at the second-lowest exit and recovered by a Turnbull NPMLE robustness. Entry is handled in the \citet{atheyseira2011} spirit: pre- and post-period bidder counts are taken as equilibrium arrival rates, and the main specification allows the post-policy SME cost distribution to differ from its pre-policy counterpart as part of the equilibrium-selection object. The exercise is conducted under observed equilibrium entry; it is not a fully solved nested entry-bidding equilibrium, and the decomposition that follows is a counterfactual-simulation accounting under the modeling assumptions stated.

Three findings organize the paper.

\noindent\textbf{(i) A \optPrefThreshold{} price preference welfare-dominates the set-aside at near-zero fiscal cost.} In thick standardized markets the simulated price shifts by $\optPrefShiftNp$ of the reference price under the preference, against $+\bneEffectTotalNp$ under the set-aside. The dominance holds under both the main equilibrium-selection and the strict-primitive-invariance treatment of the post-policy SME pool. Scaled to the post-policy run-rate of Group-65 procurement on São Paulo's R\$\,\becAnnualVolBRL~billion centralized platform, the implied annual welfare loss of the set-aside is on the order of R\$\,\welfAnnualBRLLo--\welfAnnualBRLHi~million (US\$\,\welfAnnualUSDLo--\welfAnnualUSDHi~million).

In thin heterogeneous pharmaceutical markets the comparison bifurcates. The preference dominates the set-aside under the main equilibrium-selection reading; under strict primitive invariance the ranking flips, with the implicit welfare weight required to prefer the set-aside falling to \welfWeightSiPh. The bifurcation is the diagnostic: it pins down where the modeling treatment of the protected pool is first-order, and it is exactly where the policy ranking demands caution.

\noindent\textbf{(ii) Endogenous SME entry attenuates a \latentShockRatioNp--\latentShockRatioPh~percent larger latent shock.} Hold the SME pool at its pre-policy size while imposing the SME-only rule and the simulated price effect rises to \latentShockTimesNp~percent of the realized one in non-pharma and \latentShockTimesPh~percent in pharma. Without endogenous entry the set-aside would cost \latentShockRatioNp--\latentShockRatioPh~percent more in price and welfare. Participation is the government's partial offset, not the source of the markup: the SMEs the protected segment induces to enter absorb part of the within-auction shock, though not all of it.

\noindent\textbf{(iii) The price effect operates predominantly within the auction.} Let $p_{S_1}$ be the simulated winning price under the pre-policy open regime, $p_{S_2}$ the SME-only counterfactual at the pre-policy SME pool, and $p_{S_3}$ the SME-only counterfactual at the observed post-policy SME pool. In the main specification, $p_{S_3} - p_{S_1} = +\bneEffectTotalNp$ of the reference price in non-pharma and $+\bneEffectTotalPh$ in pharma (Table~\ref{tab:v3_bne_decomp}). The within-auction piece---which I call \emph{sheltered bidding}---is \bneShareIntensNp~percent of that effect in non-pharma and \bneShareIntensPh~percent in pharma. The share stays in the \withinShareLow--\withinShareHigh~percent range across cost-distribution estimators (Turnbull NPMLE: \turnbullShareNp, \turnbullSharePh) and across the two specifications of the post-policy SME pool (strict invariance: \strictShareNp, \strictSharePh). The set-aside's price cost is mostly a property of how the surviving pool bids, not of how many firms it contains.

Three contributions follow. First, the paper provides a within-auction structural decomposition of an SME procurement preference identified from a single legal trigger, implying an annual welfare loss of R\$\,\welfAnnualBRLLo--\welfAnnualBRLHi~million on a single product group of São Paulo's R\$\,\becAnnualVolBRL~billion centralized platform. The methodological position rests on a combination of conditions rare in the procurement-preference literature: a total-exclusion legal trigger (clean separation of regimes), a Preg\~ao drop-out auction format (point-identification of cost distributions), an auction-level UH correction \citep{krasnokutskaya2011}, an observed-equilibrium-entry treatment in the \citet{atheyseira2011} spirit, and a \citet{saezstantcheva2016} welfare-weight identity. \citet{marion2007} estimates the price effect of federal SME set-asides via reduced-form variation; \citet{krasnokutskaya2011} and \citet{atheyseira2011} decompose California highway and timber bid preferences into bidding and entry channels via cross-equation constraints. This paper extends the second approach to a total-exclusion preference rule on a single clean legal trigger, allowing the welfare ranking against alternative preference instruments to be computed jointly within the same framework.

Second, the paper shows that endogenous SME entry is not a side issue but part of the welfare object: a fixed-pool counterfactual---the standard reduced-form shortcut---misstates the realized policy cost by roughly half. Third, it converts the welfare comparison into a market-design statement rather than a universal ranking. The \optPrefThreshold{} price preference welfare-dominates the set-aside in thick standardized markets across both specifications; in thin heterogeneous markets the comparison depends on how entry selection is modeled. The bifurcation identifies precisely the market conditions under which equilibrium-selection vs.\ primitive-invariance treatment of the protected pool is first-order, and converts a sensitivity check into a substantive claim about where bidder exclusion delivers proportional returns and where it does not.

The paper sits in dialogue with three lines of empirical work: recent procurement-design papers on emerging-market public contracting \citep{decarolis2024, coviello2026, ferraz2016, szerman2023}, the political economy of public contracting \citep{colonnelli2022}, and the public-finance evaluation of government programs in developing economies \citep{best2023, bandiera2021, olken2007}. The contribution here is to separate the within-auction and participation channels of an SME procurement preference on a single clean legal trigger, where the partial-preference structural literature has previously relied on cross-equation constraints to disentangle them.

A pair of comparable furosemide 40~mg purchases by the same S\~ao Paulo public buyer eight months apart and bracketing the cutoff fixes the magnitudes. In February 2018 the unit cleared roughly \furosBelowRefPct{} percent below the buyer's reference price; in October 2018 it cleared just above. The difference is one realization of the average direction the design estimates across the \nObsAuctions{} firm-auction observations and \nDistinctAuctions{} distinct auctions of the \structuralWindowM-month structural sample on each side of the \policyCutoffMonth{} cutoff---\didLowerPct{}--\didUpperPct{} percent on the absolute final price in the difference-in-differences benchmark of Appendix~\ref{app:did}, the magnitude the structural exercise of this paper interprets and prices. The headline numbers across the four reporting layers are summarized in Table~\ref{tab:preview}.

\begin{table}[!htbp]
\centering
\footnotesize
\setlength{\tabcolsep}{4pt}
\begin{threeparttable}
\caption{Headline magnitudes.}
\label{tab:preview}
\begin{tabular}{@{}p{0.65\linewidth} c c@{}}
\toprule
 & \textbf{Non-pharma} & \textbf{Pharma} \\
\midrule
DiD on $\log p^{\mathrm{final}}$ (absolute, App.~\ref{app:did}) & \multicolumn{2}{c}{$+\didLowerDec$ to $+\didUpperDec$} \\
DiD on $p^{\mathrm{final}}/p^{\mathrm{ref}}$ (App.~\ref{app:did}) & \multicolumn{2}{c}{$+\didRefDec$} \\
Simulated $p_{S_3} - p_{S_1}$, units of $p^{\mathrm{ref}}$ & $+\bneEffectTotalNp$ & $+\bneEffectTotalPh$ \\
\addlinespace
Within-auction share, main specification & \bneShareIntensNp\% & \bneShareIntensPh\% \\
Within-auction share, Turnbull NPMLE & \turnbullShareNp\% & \turnbullSharePh\% \\
Within-auction share, strict primitive invariance & \strictShareNp.0\% & \strictSharePh.0\% \\
\addlinespace
Welfare cost (\% of $p_{S_1}$, $\lambda = \lambdaMain$) & \welfLossPctLthirtyNp\% & \welfLossPctLthirtyPh\% \\
$w^{\text{SME}}_\star$ to prefer set-aside, main spec. & \welfWeightMainNp & \welfWeightMainPh \\
$w^{\text{SME}}_\star$, strict invariance & $\welfWeightSiNp$ & \welfWeightSiPh \\
\bottomrule
\end{tabular}
\begin{tablenotes}[flushleft]\footnotesize
\item DiD magnitudes from the difference-in-differences benchmark in Appendix~\ref{app:did} (against \controlGroupCount{} never-treated product groups, matched window of \structuralWindowM{} months on each side of the \policyCutoffMonth{} cutoff). Structural magnitudes from the asymmetric IPV counterfactual under observed equilibrium entry; sources by row: Sec.~\ref{sec:results} (decomposition), Sec.~\ref{sec:robustness} (Turnbull NPMLE; strict invariance), Sec.~\ref{sec:welfare} (welfare cost), Sec.~\ref{sec:discussion} (welfare weights). The pharmaceutical welfare-dominance ranking against a \optPrefThreshold{} price preference is robust under the main specification but reverses under strict primitive invariance ($w^{\text{SME}}_\star = \welfWeightSiPh < 1$); the non-pharmaceutical ranking is preserved under both specifications.
\end{tablenotes}
\end{threeparttable}
\end{table}
